\documentclass[12pt]{article}
\usepackage[utf8]{inputenc}
\usepackage[T1]{fontenc}
\usepackage{hyperref}
\title{Umbundu Language Resources for K-12 Curriculum Development}
\author{Compiled by Claude Sonnet 5 (claude-sonnet-5) for Andr\'as Kornai}
\date{}
\begin{document}
\maketitle

\section{Introduction}
Umbundu (ISO 639-3 umb) is a Bantu language of central Angola with about six million speakers. It is taught as a subject in schools. The search found grammars and Bible material but no modern school textbooks or bilingual dictionary. No search result named a native-language curriculum.

\section{Primary Educational Resources}
The search did not find Umbundu primers or school textbooks. Curriculum developers will probably have to look for them through Angolan educational authorities and publishers.

\section{Linguistic and Reference Materials}
\begin{itemize}
\item Schadeberg (1990), \emph{A Sketch of Umbundu}, an English grammatical outline of 61 pages (R\"udiger K\"oppe Verlag): \url{https://www.koeppe.de/titel_a-sketch-of-umbundu}
\item Valente (1964), \emph{Gram\'atica umbundu: a l\'ingua do centro de Angola}, a Portuguese grammar of 430 pages, listed in a library catalogue: \url{https://library.africamuseum.be/bib/264652}
\item Rosetta Project entry for Umbundu: \url{https://rosettapanglossia.longnow.org/index.php?title=Umb}
\item Language Links Database list of Umbundu resources (a historical Sanders vocabulary is listed there): \url{https://universeofmemory.com/umbundu-language-resources/}
\end{itemize}

\section{Religious and Community Resources}
\begin{itemize}
\item American Bible Society report on the new Umbundu New Testament, which the Bible Society in Angola began translating in 1999: \url{https://www.news.americanbible.org/blog/entry/god-answers-prayer-by-providing-new-testaments-to-the-umbundu-people-of-ang}
\item Umbundu Bible text (Andelale 1998 edition) at the Digital Bible Society: \url{https://dbs.org/bibles/UBUAND}. This page is titled ``Umbu Ungu'', which may be a different language (Umbu-Ungu, Papua New Guinea), so it needs checking.
\item New World Translation in Umbundu (2021): \url{https://www.jw.org/en/news/region/angola/New-World-Translation-Released-in-Umbundu/}
\end{itemize}

\section{Digital and Audio Resources}
\begin{itemize}
\item Global Recordings Network, Umbundu: \url{https://globalrecordings.net/en/language/149}
\end{itemize}

\section{Curriculum Development Context}
Umbundu is the language of the largest ethnic group in Angola and is taught as a subject. The search found no Umbundu-medium curriculum materials. Reference grammars exist in English and Portuguese, but the dictionary situation is unclear. This dossier is a starting point only.
\end{document}